\begin{helper}
Fix a finite terminal coordinate set \(U\), a nodup ordered reveal list
\(o\) with \(\operatorname{toFinset}(o)=U\), one blue support label \(B\),
and one red label \(J\).  Let \(\beta\) range over complete terminal
blue-trace labels and let \(\tau\) range over transcript labels.  For a
transcript write
\[
P_\tau,\qquad A_\tau,\qquad Z_\tau
\]
for its forced-present, forced-absent, and selected-absent coordinate sets,
and write
\[
F_{\beta,J}=B\cup \operatorname{redSupport}(\beta,J).
\]
Assume branch labels are determined by complete blue truth tables on \(U\).
Assume also the realized-cell geometry exported by the branch/transcript
construction: for every realized \((B,J,\beta,\tau)\), the red support lies
in \(U\) and is not blue, \(P_\tau,A_\tau\subseteq U\),
\(P_\tau\cap A_\tau=\emptyset\), \(F_{\beta,J}\subseteq P_\tau\), the blue
trace of \(\beta\) is encoded by \(P_\tau,A_\tau\) on every blue coordinate
of \(U\), and \(Z_\tau\subseteq A_\tau\).  Finally assume the forward
compatibility consequence: a compatible complete assignment contains
\(P_\tau\), avoids \(A_\tau\setminus Z_\tau\), and has the complete blue
trace of \(\beta\) on \(U\).

Suppose the paper's blue-first/red-second ordered reveal has supplied the
following fixed-\((B,J)\) scan interface.  There is a deterministic
transcript scan
\[
\operatorname{scan}:2^U\to\mathcal P\cup\{\bot\},
\]
four scan tests, witness assignments \(W_{\beta,\tau}\), read-dependency
sets \(D_{\beta,\tau}(e)\), prefix sets \(\operatorname{pref}(e)\), fixed
prefix supports, open-candidate sets, and a selected-present-sibling
predicate.  The prefix sets lie in \(U\).  For every realized
\((\beta,\tau)\), \(W_{\beta,\tau}\subseteq U\), every coordinate of
\(P_\tau\) is present in \(W_{\beta,\tau}\), and every coordinate of
\(A_\tau\) is absent from \(W_{\beta,\tau}\).  Every read dependency of a
coordinate is classified as one of the four ordered-reveal roles
\[
F_{\beta,J},\qquad P_\tau\setminus F_{\beta,J},\qquad
Z_\tau,\qquad A_\tau\setminus Z_\tau ,
\]
and if \(o=\operatorname{pre}\mathbin{+\!\!+}e::\operatorname{suf}\), then
\[
D_{\beta,\tau}(e)\subseteq \operatorname{toFinset}(\operatorname{pre}).
\]
Agreement with \(W_{\beta,\tau}\) on \(D_{\beta,\tau}(e)\) preserves all
four scan tests at \(e\).

Assume the scan has the two recognition properties used in the paper.  If a
residual assignment \(A_{\rm res}\subseteq U\) satisfies
\[
P_\tau\setminus F_{\beta,J}\subseteq A_{\rm res},\qquad
A_{\rm res}\cap(A_\tau\setminus Z_\tau)=\emptyset ,
\]
then
\[
\operatorname{scan}\bigl((A_{\rm res}\setminus Z_\tau)\cup F_{\beta,J}\bigr)
=\tau .
\]
If a complete assignment \(A\subseteq U\) contains \(P_\tau\), avoids
\(A_\tau\), and has the complete blue trace of \(\beta\), then
\(\operatorname{scan}(A)=\tau\).

Assume also the missing converse compatibility certificate explicitly: for
every realized \((\beta,\tau)\), if a complete assignment \(A\subseteq U\)
contains \(P_\tau\), avoids \(A_\tau\), has the complete blue trace of
\(\beta\), and the scan returns \(\tau\) on \(A\), then \(A\) is compatible
with \((B,J,\beta,\tau)\).  Finally assume the residual pruning certificate:
one residual assignment cannot reconstruct compatibly to two distinct
realized owners, and the first-mismatch interface would otherwise produce a
selected-present sibling in one of the two reconstructed assignments.

Then the residual consequences needed by the branch/transcript node hold.
Every residual assignment \(A_{\rm res}\subseteq U\) satisfying the displayed
span and avoidance conditions reconstructs to a compatible complete
assignment
\[
A^\ast=(A_{\rm res}\setminus Z_\tau)\cup F_{\beta,J}.
\]
Moreover, if one residual assignment reconstructs compatibly to two realized
pairs \((\beta,\tau)\) and \((\beta',\tau')\), then
\(\beta=\beta'\) and \(\tau=\tau'\).  The scan interface data listed above,
including the residual scan recognition, complete scan recognition, exact
read-set stability, and selected-sibling first-mismatch certificate, remain
available for downstream fixed-\((B,J)\) arguments.
\end{helper}
\begin{proof}
Fix a realized \((B,J,\beta,\tau)\) and a residual assignment
\(A_{\rm res}\subseteq U\) satisfying
\[
P_\tau\setminus F_{\beta,J}\subseteq A_{\rm res},\qquad
A_{\rm res}\cap(A_\tau\setminus Z_\tau)=\emptyset .
\]
Put
\[
A^\ast=(A_{\rm res}\setminus Z_\tau)\cup F_{\beta,J}.
\]
First \(A^\ast\subseteq U\): the residual part lies in \(U\), and the fixed
support lies in \(P_\tau\subseteq U\).  Next \(P_\tau\subseteq A^\ast\).
Indeed a coordinate of \(P_\tau\) either already lies in \(F_{\beta,J}\), in
which case it is reinserted, or lies in \(P_\tau\setminus F_{\beta,J}\), in
which case it belongs to \(A_{\rm res}\).  Such a coordinate is not selected,
because \(Z_\tau\subseteq A_\tau\) and \(P_\tau\cap A_\tau=\emptyset\).
Thus it survives the deletion of \(Z_\tau\).

The reconstructed assignment avoids all of \(A_\tau\).  If
\(e\in A_\tau\cap F_{\beta,J}\), then
\(F_{\beta,J}\subseteq P_\tau\) contradicts
\(P_\tau\cap A_\tau=\emptyset\).  If
\(e\in A_\tau\cap(A_{\rm res}\setminus Z_\tau)\), then
\(e\in A_\tau\setminus Z_\tau\) and the residual avoidance hypothesis
contradicts \(e\in A_{\rm res}\).

Now let \(e\in U\) be blue.  If \(e\) is in the blue trace of \(\beta\),
then the realized-cell geometry puts \(e\) in \(P_\tau\), hence in
\(A^\ast\).  Conversely suppose \(e\in A^\ast\).  If \(e\in F_{\beta,J}\),
then \(e\in P_\tau\), so the blue-trace equivalence gives
\(e\in\operatorname{blueTrace}(\beta)\).  If
\(e\in A_{\rm res}\setminus Z_\tau\) and \(e\) were not in the blue trace,
then the realized-cell geometry would put \(e\) in \(A_\tau\); since
\(e\notin Z_\tau\), this would put \(e\in A_\tau\setminus Z_\tau\),
contradicting residual avoidance.  Therefore \(A^\ast\) has exactly the
complete blue trace of \(\beta\).

By the residual scan-recognition clause, \(\operatorname{scan}(A^\ast)=\tau\).
The explicit converse compatibility certificate now applies to the four
facts just checked: \(A^\ast\subseteq U\), \(P_\tau\subseteq A^\ast\),
\(A^\ast\cap A_\tau=\emptyset\), and the complete blue trace.  Hence
\(A^\ast\) is compatible with \((B,J,\beta,\tau)\).

The owner-uniqueness conclusion is exactly the residual pruning certificate:
if the same residual assignment reconstructed compatibly to two realized
pairs, the selected-sibling first-mismatch alternative would be pruned by
the fixed-prefix support data, so the two owners must be the same.  The
remaining displayed scan/read properties are part of the supplied ordered
reveal interface.  The dependencies used here are the scan construction and
read stability from
\(\noderef{FixedSetHistoryCellRISIRevealOrderScanConstruction}\) and
\(\noderef{FixedSetHistoryCellRISIResidualScanStability}\), the fixed
red-branch scan interface
\(\noderef{FixedSetHistoryCellRedBranchScanInterfaceConstruction}\), the
selected-sibling pruning interface
\(\noderef{FixedSetHistoryCellRedPrefixScanInvariant}\) and
\(\noderef{FixedSetHistoryCellFixedBJResidualOverlapConsistency}\), the
deterministic residual scan-output package
\(\noderef{FixedSetHistoryCellRedResidualDeterministicScanOutputs}\), and
the terminal assignment extension
\(\noderef{FixedSetHistoryCellTerminalAssignmentExtension}\).
\end{proof}
